% To be reformatted in the journal style upon acceptance.
\documentclass[11pt]{amsart}

\usepackage{amsmath,amssymb}
\usepackage[hidelinks]{hyperref}

\newtheorem{thm}{Theorem}[section]
\newtheorem{lem}[thm]{Lemma}
\newtheorem{cor}[thm]{Corollary}
\newtheorem{prop}[thm]{Proposition}
\theoremstyle{definition}
\newtheorem{defi}[thm]{Definition}
\newtheorem{rem}[thm]{Remark}
\newtheorem{ques}[thm]{Question}
\newtheorem*{thmA}{Theorem A}
\newtheorem*{thmB}{Theorem B}

\newcommand{\I}{\mathbb{I}}
\newcommand{\DM}{\mathrm{DM}}
\newcommand{\KL}{\mathrm{KL}}
\newcommand{\BA}{\mathrm{BA}}
\newcommand{\comp}{\mathsf{comp}}
\newcommand{\inc}{\mathsf{inc}}
\newcommand{\Ldash}{L}

\begin{document}

\title[The reversal quotient is a $K(\mathbb{Z}/2,1)$]{Cubical
  models with reversal do not present spaces:\\ the interval
  quotient is a $K(\mathbb{Z}/2, 1)$}

\author{Ryota Shibusawa}
\address{Daiichi Institute of Technology, Japan}
\email{r-shibusawa@daiichi-koudai.ac.jp}

\subjclass[2020]{55U35; 18N40, 55U10, 06D30, 03B38}
\keywords{cubical sets, cubical type theory, model structure,
  reversal, De Morgan algebra, Kleene algebra, geometric
  realization, Eilenberg--MacLane space, test category}

\begin{abstract}
  Since 2018 it has been folklore, recorded only in a lecture and a
  mailing-list discussion, that the cubical-type model structure on
  De Morgan cubical sets --- the semantic home of CCHM cubical type
  theory --- does not present the homotopy theory of spaces: the
  quotient $\Ldash$ of the representable interval by the reversal
  $x \mapsto \neg x$ has contractible geometric realization but is
  not weakly contractible.  We give the first complete proof, in a
  form that covers the De Morgan, Kleene and Boolean cube categories
  uniformly --- the Kleene case having never been examined --- and
  we identify the homotopy type: $\Ldash$ is a
  $K(\mathbb{Z}/2, 1)$ in the cubical-type model structure.  The
  proof has two engines.  A one-line \emph{freeness theorem} shows
  the reversal acts without fixed cells on the interval, in every
  dimension and all three theories; a finite \emph{reachability
  census} of reversal-symmetric squares then closes the
  non-contractibility argument, adapting Coquand's
  trivial-cofibration/fibration factorization method from the
  cartesian symmetric square to the symmetric line --- including two
  subtleties absent from the cartesian case (connection-degenerate
  composition directions, and symmetric contractibility of the
  connection-degenerate line, which makes the reachability census
  essential).  Freeness then powers a \emph{sheet-transport}
  fibration structure on the covering $\square^1 \to \Ldash$, whose
  total space is contractible by a connection and whose fiber is the
  discrete $\mathbb{Z}/2$: the quotient computes the homotopy
  quotient, not the orbit space.  Consequences: the cubical-type and
  test model structures differ on all three cube categories, and
  geometric realization is not a Quillen equivalence for any of
  them.
\end{abstract}

\maketitle

\section{Introduction}
\label{sec:intro}

Cubical type theory in the style of
Cohen--Coquand--Huber--M\"ortberg \cite{CCHM} is interpreted in
presheaves on a cube category whose interval carries connections
$\wedge, \vee$ and a reversal $\neg$ --- the free De Morgan algebra
structure.  These presheaves carry a \emph{cubical-type model
structure} (Sattler \cite{SattlerEEP}; Coquand \cite{CoqNote} gives
an explicit trivial-cofibration/fibration factorization), whose
fibrant objects model univalent type theory.  Whether this model
structure presents the homotopy theory of spaces was answered in
the negative, for the cartesian cube category, by an argument of
Sattler and Coquand in June 2018: the quotient of the square by the
coordinate swap is not weakly contractible, though its realization
is a triangle \cite{CoqNote, CoqML}.  The equivariant repair of
this defect for cartesian cubes is the recent landmark of
Awodey--Cavallo--Coquand--Riehl--Sattler \cite{ACCRS}.

For the De Morgan cube category the corresponding negative claim
--- with counterexample the quotient
\[
  \Ldash \;:=\; \square^1 / (x \sim \neg x)
\]
of the representable interval by the reversal --- has circulated
since the same discussions \cite{CoqML, SatTalk}, but no proof has
been published: Coquand's note proves the cartesian case in full
and says of the De Morgan case only that the argument ``should also
be valid'' \cite[Rem.~3]{CoqNote}; the recent literature
\cite{ACCRS, CavalloSattlerRev} continues to cite the 2018 lecture
\cite{SatTalk} for it.  The Kleene cube category --- the canonical
cube category of the standard interval $[0,1]$
\cite{BuchholtzMorehouse} --- appears in this connection only once
in the literature, as a parenthetical remark in the 2018 thread
\cite{CoqML}.

This paper closes the question, in a stronger form than the
folklore records.  Throughout, $\mathcal{T}$ denotes any of the De
Morgan, Kleene, or Boolean cube categories, and ``the model
structure'' means the cubical-type model structure on presheaves
over $\mathcal{T}$.

\begin{thmA}[= Theorem~\ref{thm:nogo}]
  In all three theories, $\Ldash \to 1$ is not a weak equivalence,
  although the geometric realization $|\Ldash| \cong [0, 1/2]$ is
  contractible.  Hence the cubical-type model structure differs
  from the test model structure, and geometric realization is not a
  Quillen equivalence.
\end{thmA}

\begin{thmB}[= Theorem~\ref{thm:kz2}]
  In all three theories, the quotient map
  $p : \square^1 \to \Ldash$ is a fibration with contractible total
  space and fiber the discrete $\mathbb{Z}/2$; consequently
  $\Ldash$ is a $K(\mathbb{Z}/2, 1)$ in the cubical-type model
  structure.
\end{thmB}

Theorem A for De Morgan and Boolean cubes is the first published
proof of the folklore claim; for Kleene cubes it is new.  Theorem B
identifies the actual homotopy type, which no source has claimed:
the moral is that \emph{the presheaf quotient computes the homotopy
quotient of the reversal action, not its orbit space} --- exactly
the phenomenon that equivariance was invented to control on the
cartesian side \cite{ACCRS}, here appearing for the
$\mathbb{Z}/2$-action of the reversal, for which no equivariant
theory yet exists.

\subparagraph*{Method.}  Everything flows from one observation
(Theorem~\ref{thm:free}): \emph{the reversal acts freely on the
cells of the interval} --- no reparametrization $t$ of the interval
satisfies $t = \neg t$, in any dimension and any of the three
theories, because the all-consistent vertex assignment
is fixed by the point-level involution while Boolean values are
not.  Freeness has two independent uses.  For Theorem A we adapt
Coquand's invariant method \cite[\S4--5]{CoqNote} from symmetric
squares to \emph{symmetric lines}; two subtleties absent from the
cartesian case appear.  First, the composition directions fixed by
the reversal now include the connection-degenerate elements
$x \wedge \neg x$ and $x \vee \neg x$, enlarging the case analysis
of the factorization argument.  Second --- and this is where the
cartesian intuition genuinely fails --- the connection-degenerate
line $\ell(x \wedge \neg x)$ \emph{is} symmetrically contractible,
so the chain version of the invariant could conceivably connect
the generator to a constant through it; whether it can is a finite
question about two-variable free algebras, and an exhaustive
\emph{reachability census} answers it: the symmetric-square graph
on symmetric lines decomposes, in all three theories, into the
isolated component of the generator and the contractible component
of the degenerate lines.  (The census sizes are small: $60$, $32$
and $8$ symmetric squares for De Morgan, Kleene, Boolean.)  For
Theorem B, freeness makes the two lifts of any cell of $\Ldash$
into disjoint \emph{sheets}, and the fibration structure on $p$ is
\emph{sheet transport}: the unique lift through a prescribed face,
with a tube-forcing lemma showing all compatibility data lies on
the same sheet.  The reachability census and the monodromy of the
covering tell one story: the generator line has odd parity, the
degenerate lines even --- the combinatorial shadow of
$\pi_1 = \mathbb{Z}/2$.

\subparagraph*{Relation to the literature.}  Our debt to
\cite{CoqNote} is total as to method; the contribution is the
adaptation (with its two new subtleties), the uniform treatment of
the three theories, the previously unexamined Kleene case, the
identification of the homotopy type, and machine verification of
every finite ingredient.  The reversal-elimination theorem of
Cavallo--Sattler \cite{CavalloSattlerRev} obtains a model of a
theory \emph{with} reversals presenting spaces by interpreting over
reversal-free cubes; Theorems A and B delimit what it is they are
eliminating: reversal as a \emph{cube-category} symmetry inevitably
manufactures $K(\mathbb{Z}/2,1)$s from free quotients.  Whether an
equivariant fibration structure for the reversal action --- the
$\mathbb{Z}/2$, or in higher dimensions hyperoctahedral, analogue
of \cite{ACCRS} --- repairs the De Morgan and Kleene models is, we
believe, the natural open problem these results isolate
(\S\ref{sec:outlook}).  The literal-cube combinatorics of free De
Morgan and Kleene algebras used throughout is developed in
\cite{DMFibers, Kleene}.

\section{Preliminaries}
\label{sec:prelim}

\subsection{Cube categories and the interval}

Following \cite[\S1]{CoqNote}, the base category has as objects
finite sets of names $I, J, \dots$; a morphism $f : J \to I$ is a
function from $I$ to the free algebra $\mathcal{T}(J)$ of the
interval theory $\mathcal{T}$ on the names $J$, where
$\mathcal{T}(J)$ is the free De Morgan algebra $\DM(J)$, the free
Kleene algebra $\KL(J)$, or the free Boolean algebra $\BA(J)$.  The
formal interval is $\I(J) = \mathcal{T}(J)$.  We use the
literal-cube representations \cite{DMFibers, Kleene}: $\DM(J)$ is
the lattice of monotone Boolean functions on $\{0,1\}^{2|J|}$
(coordinates in pairs $(x, \neg x)$), $\KL(J)$ its quotient by the
mixed points, $\BA(J)$ the functions on consistent points.  The
reversal $\neg$ acts on elements by
$(\neg t)(\alpha) = 1 - t(\sigma\bar\alpha)$, where $\bar\alpha$
complements all coordinates and $\sigma$ swaps each pair.

$\square^1 = \mathrm{Yon}(x)$ denotes the representable interval,
with cells $\square^1(J) = \mathcal{T}(J)$, and
\[
  \Ldash \;=\; \square^1 / (t \sim \neg t),
  \qquad \Ldash(J) = \mathcal{T}(J)/\mathord{\sim},
\]
the quotient presheaf (colimits are computed pointwise).  We write
$[t]$ for cells of $\Ldash$ and $\ell(t)$ when we think of them as
reparametrizations of the generating line.  Note
$[0] = [1]$: $\Ldash$ has a single vertex $\ell_0$, and the
generator $[x]$ is a loop at it.

\subsection{The model structure and the factorization}

The cubical-type model structure on
$\widehat{\mathcal{T}}$ has cofibrations the (decidable)
monomorphisms; its fibrations are defined by a uniform composition
operation $\comp^{k \to l}$ against partial elements, and Coquand's
note constructs, for any map $\sigma : A \to B$, an explicit trivial
cofibration--fibration factorization $A \to T \to B$ where $T$ is
the total space of an inductively defined family $E$ over $B$ with
constructors $\inc$ and $\comp^{k \to l}(u_k, \psi, u, v)$
($k \neq l \in \I(J)$, $\psi \neq 1$), and restriction maps given
by: reindexing when $kg \neq lg$ and $\psi g \neq 1$;
$u_k g$ when $kg = lg$; $u_{[l]g}$ when $\psi g = 1$
\cite[\S3]{CoqNote}.  A map is a weak equivalence iff in (any) such
factorization the fibration part is a trivial fibration; trivial
fibrations have sections.  All of this is verbatim uniform in
$\mathcal{T}$: the constructions never inspect the interval theory
beyond its algebra operations.

\section{The freeness theorem}
\label{sec:free}

\begin{thm}[freeness]
  \label{thm:free}
  In all three theories and for every $J$, no $t \in \mathcal{T}(J)$
  satisfies $t = \neg t$.  Hence the reversal action on the cells of
  $\square^1$ is free, and every cell of $\Ldash$ has exactly two
  lifts $\{t, \neg t\}$.
\end{thm}

\begin{proof}
  Let $\alpha_0$ be the point of the literal cube with every pair
  equal to $(0,1)$ --- the consistent point of the all-zeros
  assignment, present in all three representations.  Then
  $\sigma\bar\alpha_0 = \alpha_0$: complementing gives all pairs
  $(1,0)$, and swapping each pair returns $(0,1)$.  So
  $t = \neg t$ would give
  $t(\alpha_0) = 1 - t(\sigma\bar\alpha_0) = 1 - t(\alpha_0)$,
  impossible for a Boolean value.
\end{proof}

(Machine census: there are $0$ self-dual elements among all
$7{,}828{,}354$ elements of $\DM(3)$, all $43{,}918$ of $\KL(3)$,
and below; see \S\ref{sec:machine}.)

\section{Theorem A: $\Ldash$ is not weakly contractible}
\label{sec:nogo}

\begin{defi}
  A square $c(x, z)$ \emph{symmetrically connects} lines $a(x)$ and
  $b(x)$ if $c(x, 0) = a$, $c(x,1) = b$ (or vice versa) and
  $c(\neg x, z) = c(x, z)$.  Write $P'(A)$ for the property: every
  reversal-symmetric line $a(x) = a(\neg x)$ of $A$ is connected to
  a constant line by a finite chain of symmetrically connecting
  squares.  (In a quotient like $\Ldash$, symmetry of $[t]$ means
  $t \circ \sigma_x \in \{t, \neg t\}$.)
\end{defi}

\begin{lem}
  \label{lem:section}
  If $u : A \to B$ has a section then $P'(A)$ implies $P'(B)$.
\end{lem}

\begin{proof}
  As in \cite[Lem.~4.1]{CoqNote}: push the symmetric line along the
  section, connect there, map the chain back.
\end{proof}

\begin{prop}
  \label{prop:pres}
  For every $\sigma : A \to B$ with factorization $A \to T \to B$
  as above, $P'(A)$ implies $P'(T)$.
\end{prop}

\begin{proof}
  A symmetric line of $T$ over $v[l]$ is either
  $(\sigma a, \inc\, a)$ with $a$ a symmetric line of $A$ ---
  handled by $P'(A)$ --- or
  $(v(x, l), \comp^{k \to l}(a, \psi, u, \langle z\rangle v(x,z)))$
  with $k \neq l \in \I(x)$.  Two remarks before the induction.
  First, the substitution $g = (x \mapsto \neg x)$ is an
  \emph{automorphism} of $\mathcal{T}(J)$, so the degenerate
  restriction clauses cannot fire under it: $kg = lg$ iff $k = l$,
  and $\psi g = 1$ iff $\psi = 1$.  Symmetry of the cell is
  therefore an equality of $\comp$-constructors, forcing
  $k \circ \sigma_x = k$, $l \circ \sigma_x = l$, $a$ symmetric,
  $\psi$, $u$, $v$ symmetric.  Second, the fixed elements of
  $\I(x)$ under $x \mapsto \neg x$ are $\{0, 1\}$ in the Boolean
  theory but
  \[
    \{\,0,\ 1,\ x \wedge \neg x,\ x \vee \neg x\,\}
  \]
  in the De Morgan and Kleene theories: the composition may enter
  or exit through a connection-degenerate direction, a case with no
  cartesian analogue.  The induction is nevertheless uniform: with
  a fresh name $t$, the square
  \[
    c(x, t) \;=\;
    \bigl(v(x,t),\ \comp^{k \to t}(a, \psi, u,
    \langle z\rangle v(x,z))\bigr)
  \]
  is symmetric (all its data is), equals the given line at $t = l$,
  and at $t = k$ equals $(v[k], a)$ by the $kg = lg$ restriction
  clause.  Since $a$ has smaller inductive height and is symmetric,
  induction reaches an $\inc$-cell and $P'(A)$ concludes.
\end{proof}

\begin{prop}
  \label{prop:census}
  $P'(\Ldash)$ fails, in all three theories.
\end{prop}

\begin{proof}
  The symmetric lines of $\Ldash$ are exactly its three (Boolean:
  two) lines: the generator $[x]$, the connection-degenerate
  $[x \wedge \neg x] = [x \vee \neg x]$ (absent in the Boolean
  case), and the constant $[0] = [1]$: every element of
  $\mathcal{T}(x)$ is orbit-symmetric.  A symmetrically connecting
  square of $\Ldash$ is an orbit $[t]$, $t \in \mathcal{T}(x,z)$,
  with $t \circ \sigma_x \in \{t, \neg t\}$ --- a \emph{finite} set,
  and the faces $z = 0, 1$ of its members are computable.  The
  exhaustive census (\S\ref{sec:machine}): the symmetric squares
  number $60$ (De Morgan), $32$ (Kleene), $8$ (Boolean), and the
  induced graph on symmetric lines is, in each theory,
  \[
    \underbrace{[x] \circlearrowleft}_{\text{isolated}}
    \qquad\qquad
    [x \wedge \neg x] \longleftrightarrow \text{const}
    \,\circlearrowleft ,
  \]
  i.e.\ every symmetric square with one face $[x]$ has both faces
  $[x]$.  Two features deserve note.  The degenerate line
  \emph{is} symmetrically contractible ---
  $t = (x \wedge \neg x) \wedge \neg z$ is exactly symmetric and
  connects it to $0$ --- so the isolation of $[x]$ is essential
  and not a foregone conclusion; and the connection trick that
  kills the cartesian symmetric-square invariant
  \cite[Rem.~1]{CoqNote} produces $\ell(x \vee z)$, which is
  \emph{not} symmetric ($x \vee z \neq \neg x \vee z$ and
  $\neq \neg(x \vee z)$ in the free algebras), in accordance with
  the census.  Hence no chain connects $[x]$ to a constant.
\end{proof}

\begin{thm}
  \label{thm:nogo}
  In all three theories $\Ldash \to 1$ is not a weak equivalence;
  the cubical-type model structure differs from the test model
  structure; and geometric realization is not a Quillen
  equivalence.
\end{thm}

\begin{proof}
  If $\Ldash \to 1$ were a weak equivalence then so would be the
  vertex $1 \to \Ldash$ (two-out-of-three); factor it as
  $1 \to T \to \Ldash$; the fibration $T \to \Ldash$ is then
  trivial, hence has a section.  $P'(1)$ holds trivially,
  Proposition~\ref{prop:pres} gives $P'(T)$, and
  Lemma~\ref{lem:section} gives $P'(\Ldash)$ --- contradicting
  Proposition~\ref{prop:census}.  For the comparisons:
  $|\Ldash| \cong [0,1]/(s \sim 1-s) \cong [0, 1/2]$ is
  contractible, and all three cube categories are strict test
  categories \cite{BuchholtzMorehouse}, so in the test model
  structure $\Ldash \to 1$ \emph{is} a weak equivalence --- the two
  localizations differ.  Since cofibrations are monomorphisms,
  $\Ldash$ is cofibrant, and a left Quillen equivalence reflects
  weak contractibility of cofibrant objects
  \cite[1.3.16]{Hovey}; realization fails to do so.
\end{proof}

\section{Theorem B: $\Ldash$ is a $K(\mathbb{Z}/2, 1)$}
\label{sec:kz2}

\begin{thm}[sheet transport]
  \label{thm:fib}
  The quotient $p : \square^1 \to \Ldash$ is a fibration: it
  carries a uniform composition structure.
\end{thm}

\begin{proof}
  In dependent form, the fiber over a cell $v \in \Ldash(J)$ is
  $E(v) = \{t, \neg t\}$, the two lifts (Theorem~\ref{thm:free}).
  Given composition data --- $v \in \Ldash(J^{+})$,
  $k \neq l \in \I(J)$, a lift $u_k \in E(v[k])$, $\psi \neq 1$,
  and a compatible partial family $u_f \in E(vf)$ for
  $f : K \to J^{+}$ with $\psi \mathsf{p} f = 1$ --- there is a
  \emph{unique} lift $T \in \mathcal{T}(J^{+})$ of $v$ with
  $T[k] = u_k$: of the two lifts $T, \neg T$ exactly one restricts
  to $u_k$, since $T[k] = \neg T[k]$ would contradict freeness at
  level $J$.  Define
  \[
    \comp^{k \to l}(u_k, \psi, u) \;:=\; T[l].
  \]
  \emph{Tube forcing}: for $g : K \to J$ with $\psi g = 1$, the
  tube cell $u_{g^{+}}$ (part of the data, as
  $\psi \mathsf{p} g^{+} = 1$) satisfies
  $u_{g^{+}}[kg] = u_{[k]g} = u_k g = T[k] g = (T g^{+})[kg]$,
  using $[k] \circ g = g^{+} \circ [kg]$; by freeness at level
  $K^{+}$, agreeing with $T g^{+}$ at one face forces
  $u_{g^{+}} = T g^{+}$, whence
  $u_{[l]g} = u_{g^{+}}[lg] = (T g^{+})[lg] = T[l] g$: the
  composition agrees with the prescribed data on $\psi$.
  Uniformity: for $g$ with $kg \neq lg$, $\psi g \neq 1$, the
  unique lift of $v g^{+}$ through $u_k g$ is $T g^{+}$, so
  $\comp(\dots) g = T[l] g = \comp(u_k g, \psi g, u g^{+})$ ---
  naturality of the construction.  Finally
  $\comp^{k \to k} = T[k] = u_k$.
\end{proof}

\begin{thm}
  \label{thm:kz2}
  In all three theories, $\Ldash$ is a $K(\mathbb{Z}/2, 1)$: the
  loop object of its fibrant replacement at the vertex is weakly
  equivalent to the discrete two-point object, and its higher
  homotopy vanishes.
\end{thm}

\begin{proof}
  $p$ is a fibration (Theorem~\ref{thm:fib}) with fiber over the
  vertex the discrete two-point presheaf.  Its total space
  $\square^1$ is contractible: the connection square
  $c(x, z) = x \wedge \neg z$ is a homotopy from the identity to
  the constant $0$.  The fiber sequence
  $\mathbb{Z}/2 \to \square^1 \to \Ldash$ then identifies, in the
  homotopy category of the model structure, the loop object of
  $\Ldash$ at $\ell_0$ with the discrete fiber and kills the
  higher homotopy: $\pi_n(\Ldash) = \pi_{n-1}(\mathbb{Z}/2) = 0$
  for $n \geq 2$.  Concretely, the generator loop $[x]$ lifts to
  the path $x$ from $0$ to $1$, crossing sheets: its monodromy is
  the nontrivial element of $\mathbb{Z}/2$ (this re-proves
  Theorem~\ref{thm:nogo} independently of \S\ref{sec:nogo}), while
  $[x] \cdot [x]$ lifts to a genuine loop and is
  null-homotopic --- $\pi_1 = \mathbb{Z}/2$ exactly.
\end{proof}

\begin{rem}
  The two proofs of non-contractibility are one: the reachability
  census of Proposition~\ref{prop:census} separates the symmetric
  lines precisely by the parity of their monodromy --- $[x]$ lifts
  $0 \to 1$ (odd), while $[x \wedge \neg x]$ and the constants lift
  to loops (even) --- and symmetric squares preserve the parity.
  The census is the combinatorial shadow of
  $\pi_1 = \mathbb{Z}/2$.
\end{rem}

\begin{rem}
  Theorem~\ref{thm:kz2} says the presheaf quotient of the interval
  by the reversal computes the \emph{homotopy} quotient
  $(\square^1)_{h\mathbb{Z}/2} \simeq B\mathbb{Z}/2$ rather than
  the orbit space --- because the action is free on cells
  (Theorem~\ref{thm:free}) even though its topological realization
  has a fixed point at $\tfrac12$.  This is precisely the mechanism
  that equivariant fibrations neutralize on the cartesian side
  \cite{ACCRS}, where the symmetric quotients $I^n/H$ play the role
  of $\Ldash$; the reversal case is, if anything, cleaner, since
  freeness makes the covering honest.  No equivariant theory for
  the reversal ($\mathbb{Z}/2$ per direction; hyperoctahedrally,
  $\Sigma_n \ltimes (\mathbb{Z}/2)^n$ in dimension $n$) exists in
  the literature.
\end{rem}

\section{Machine verification}
\label{sec:machine}

All finite ingredients were verified exhaustively (scripts
accompany the artifact).  \emph{Freeness}: $0$ self-dual elements
in $\DM(m)$ for $m \leq 3$ (out of $6$, $168$, $7{,}828{,}354$) and
in $\KL(m)$ for $m \leq 3$ (out of $6$, $84$, $43{,}918$) ---
the theorem's one-line proof makes higher $m$ unnecessary, but the
census independently guards the representations.  \emph{No rescue
terms}: no $t$ in $\DM(x,z)$ or $\KL(x,z)$ is
reversal-compatible, restricts to the generator at $z = 0$, and is
constant at $z = 1$ --- the connection trick of
\cite[Rem.~1]{CoqNote} has no symmetric analogue.
\emph{Reachability}: the symmetric squares number $60$ / $32$ / $8$
(De Morgan / Kleene / Boolean); their face graphs on symmetric
lines have edge counts $[x] \to [x]$: $10$/$6$/$4$, degenerate
$\leftrightarrow$ constant: $10{+}10$ / $6{+}6$ / --- , constant
$\to$ constant: $10$/$8$/$4$, and \emph{no} edge leaves the
$[x]$-component in any theory.  \emph{L-cells}: the squares of
$\Ldash$ over the generator number $62$ ($44$ with $z$-varying
faces, $12$ contracting the generator non-symmetrically) --- the
quotient is rich; only symmetry starves it.

\section{Outlook}
\label{sec:outlook}

Three directions.  (1) \emph{Reversal equivariance}: define uniform
fibrations equivariant for the $\mathbb{Z}/2$-actions of reversal
(hyperoctahedrally, $\Sigma_n \ltimes (\mathbb{Z}/2)^n$) on De
Morgan or Kleene cubes, in the style of \cite{ACCRS}; the freeness
theorem says the reversal orbits are regular on the interval, so
the equivariant filling conditions have no isotropy at the crucial
cells --- we regard the repair as plausible and the natural
continuation.  (2) \emph{The remaining open positive case}: our
argument uses reversal essentially; for Dedekind cubes (no
reversal) the presentation question remains open
\cite{StreicherWeinberger, CavalloSattlerElegance}, with a
constructive solution announced by Sattler.  (3) \emph{Stronger
inequivalence}: Theorems A and B show the canonical comparison
fails; whether the cubical-type localization is inequivalent to
spaces under \emph{every} functor --- asserted in the folklore, and
explicitly not excluded on the nLab --- would follow from an
invariant distinguishing the localizations; the abundance of
$K(\mathbb{Z}/2,1)$s with contractible realization produced here
(one for every free reversal quotient) is a start.

\subparagraph*{Artifact.}  The verification scripts and the
machine-checked companions are available in the artifact
repository:\\
\url{https://github.com/r-shibusawa/cubical-strict-j}, archived
at DOI \href{https://doi.org/10.5281/zenodo.21405961}%
{10.5281/zenodo.21405961}.

\begin{thebibliography}{10}

\bibitem{ACCRS}
S.~Awodey, E.~Cavallo, T.~Coquand, E.~Riehl, C.~Sattler.
\newblock The equivariant model structure on cartesian cubical
  sets.
\newblock \emph{Adv.\ Math.}\ 495 (2026), 110965.

\bibitem{BuchholtzMorehouse}
U.~Buchholtz, E.~Morehouse.
\newblock Varieties of cubical sets.
\newblock \emph{RAMiCS 2017}, LNCS 10226, 2017.

\bibitem{CavalloSattlerElegance}
E.~Cavallo, C.~Sattler.
\newblock Relative elegance and cartesian cubes with one
  connection.
\newblock \emph{Canad.\ J.\ Math.}, to appear; arXiv:2211.14801.

\bibitem{CavalloSattlerRev}
E.~Cavallo, C.~Sattler.
\newblock Eliminating reversals from cubical type theories.
\newblock \emph{LICS 2026}, LIPIcs 380, 2026.

\bibitem{CCHM}
C.~Cohen, T.~Coquand, S.~Huber, A.~M\"ortberg.
\newblock Cubical type theory: a constructive interpretation of the
  univalence axiom.
\newblock \emph{TYPES 2015}, LIPIcs 69, 2018.

\bibitem{CoqNote}
T.~Coquand.
\newblock Trivial cofibration-fibration factorization with one
  application.
\newblock Unpublished note, June 2018; circulated on the Homotopy
  Type Theory mailing list \cite{CoqML}.

\bibitem{CoqML}
T.~Coquand et al.
\newblock Quillen model structure.
\newblock Homotopy Type Theory mailing-list discussion, June 2018.
\newblock \url{https://groups.google.com/g/homotopytypetheory/c/RQkLWZ_83kQ}.

\bibitem{Hovey}
M.~Hovey.
\newblock \emph{Model Categories}.
\newblock Math.\ Surveys Monogr.\ 63, AMS, 1999.

\bibitem{SatTalk}
C.~Sattler.
\newblock Do cubical models of type theory also model homotopy
  types?
\newblock Lecture, Hausdorff Trimester Program ``Types, Sets and
  Constructions'', Bonn, June 2018.

\bibitem{SattlerEEP}
C.~Sattler.
\newblock The equivalence extension property and model structures.
\newblock Preprint, arXiv:1704.06911, 2017.

\bibitem{StreicherWeinberger}
T.~Streicher, J.~Weinberger.
\newblock Simplicial sets inside cubical sets.
\newblock \emph{Theory Appl.\ Categ.}\ 37 (2021), 276--286.

\bibitem{DMFibers}
R.~Shibusawa.
\newblock Boundary fibers in free De Morgan algebras: Boolean
  intervals, Sperner bounds, and median contractions.
\newblock Preprint, 2026.  In the artifact repository (release
  \texttt{v1.6.0}, DOI
  \href{https://doi.org/10.5281/zenodo.21695307}%
  {10.5281/zenodo.21695307}).

\bibitem{Kleene}
R.~Shibusawa.
\newblock The Kleene interpolation: boundary fibers and the strict
  layer across the subvarieties of De Morgan algebras.
\newblock Preprint, 2026.  In the artifact repository (release
  \texttt{v1.8.1}, DOI
  \href{https://doi.org/10.5281/zenodo.21695917}%
  {10.5281/zenodo.21695917}).

\end{thebibliography}

\end{document}
